\chapter{相关工作}\label{chap:related}

\section{期权定价模型}\label{sec:option_models}

\subsection{Black-Scholes-Merton模型}

Black、Scholes与Merton于1973年提出了期权定价的奠基性框架\cite{black1973pricing,merton1973theory}。假设标的资产价格$S$服从几何布朗运动：
\begin{equation}
  dS = \mu S\,dt + \sigma S\,dW_t,
\end{equation}
其中$\sigma$为常数波动率，$W_t$为标准布朗运动。通过无套利原理，欧式期权价值$V(S,t)$满足BSM方程：
\begin{equation}\label{eq:bsm}
  \frac{\partial V}{\partial t}
  + \frac{1}{2}\sigma^2 S^2 \frac{\partial^2 V}{\partial S^2}
  + rS\frac{\partial V}{\partial S}
  - rV = 0,
\end{equation}
其中$r$为无风险利率。欧式看涨期权具有封闭解析解（Black-Scholes公式）：
\begin{equation}\label{eq:bs_formula}
  V_\text{BS}(S,t) = S\Phi(d_1) - Ke^{-r(T-t)}\Phi(d_2),
\end{equation}
其中$d_1 = \frac{\ln(S/K)+(r+\sigma^2/2)(T-t)}{\sigma\sqrt{T-t}}$，$d_2=d_1-\sigma\sqrt{T-t}$，$\Phi(\cdot)$为标准正态累积分布函数。BSM模型的主要局限在于假设波动率为常数，无法解释市场中普遍存在的波动率微笑现象。

\subsection{常弹性方差模型}

Cox和Ross\cite{cox1976valuation}提出了常弹性方差（CEV）模型，将扩散系数推广为依赖于标的价格的局部波动率：
\begin{equation}
  dS = \mu S\,dt + \sigma S^{\beta}\,dW_t,
\end{equation}
其中$\beta$为弹性参数。当$\beta=1$时退化为BSM；当$\beta<1$时，波动率随股价上涨而下降，能够捕捉股票市场的杠杆效应。对应的期权定价PDE为：
\begin{equation}\label{eq:cev}
  \frac{\partial V}{\partial t}
  + \frac{1}{2}\sigma^2 S^{2\beta} \frac{\partial^2 V}{\partial S^2}
  + rS\frac{\partial V}{\partial S}
  - rV = 0.
\end{equation}
本文取$\beta=0.5$（平方根过程）。Schroder\cite{schroder1989computing}指出，$\beta=0.5$时CEV欧式看涨期权存在精确解析解（非中心卡方分布），本文以此作为CEV的参考基准。

\subsection{Heston随机波动率模型}

Heston\cite{heston1993closed}将波动率建模为均值回归过程（CIR过程）：
\begin{align}
  dS &= \mu S\,dt + \sqrt{v}\,S\,dW_t^S, \\
  dv &= \kappa(\theta - v)\,dt + \xi\sqrt{v}\,dW_t^v,
\end{align}
其中$v$为瞬时方差，$\kappa$为均值回归速度，$\theta$为长期均值，$\xi$为vol-of-vol，$dW_t^S$与$dW_t^v$的相关系数为$\rho$。期权价值$V(S,v,t)$满足二维PDE：
\begin{equation}\label{eq:heston}
  \frac{\partial V}{\partial t}
  + \frac{1}{2}vS^2\frac{\partial^2 V}{\partial S^2}
  + \rho\xi v S\frac{\partial^2 V}{\partial S\partial v}
  + \frac{1}{2}\xi^2 v\frac{\partial^2 V}{\partial v^2}
  + rS\frac{\partial V}{\partial S}
  + \kappa(\theta-v)\frac{\partial V}{\partial v}
  - rV = 0.
\end{equation}
Heston模型具有半解析解（特征函数方法与Gil-Pelaez反演），能够产生波动率微笑和厚尾收益率分布，在指数期权和外汇期权中被广泛使用。Feller条件$2\kappa\theta>\xi^2$保证方差过程$v$始终为正。

\section{物理信息神经网络}\label{sec:pinn}

\subsection{基本原理}

物理信息神经网络（PINN）由Raissi等人\cite{raissi2019physics}于2019年提出，其核心思想是将PDE残差直接嵌入神经网络的损失函数。对于PDE $\mathcal{F}[u]=0$，PINN的总损失为：
\begin{equation}\label{eq:pinn_loss}
  \mathcal{L} = w_{\text{pde}}\mathcal{L}_{\text{pde}}
              + w_{\text{ic}}\mathcal{L}_{\text{ic}}
              + w_{\text{bc}}\mathcal{L}_{\text{bc}},
\end{equation}
其中各项分别为配点处的PDE残差、初值条件和边界条件的均方误差，$w_{\text{pde}},w_{\text{ic}},w_{\text{bc}}$为权重超参数。PDE残差通过自动微分精确计算，无需有限差分近似。PINN无需网格剖分，训练完成后对任意输入的推断时间约为毫秒级。

\subsection{PINN在期权定价中的应用}

Dhiman和Hu\cite{dhiman2023pinn}率先将PINN应用于欧式和美式期权定价，提出了门控网络（gated network）结构，通过浅层分支与深层分支的加权组合分别捕捉解的简单特征与复杂特征。Zamxaka\cite{zamxaka2023cev}系统研究了PINN对BSM和CEV模型的求解，构建了参数化PINN，使单个网络可在不同参数设置下直接推断期权价格。Bae等人\cite{bae2024localvol}将PINN扩展至BSM、CEV及Dupire局部波动率模型，实现了期权价格与Greeks的同步计算。针对Heston模型，Hainaut和Casas\cite{hainaut2024heston}构建了参数化PINN，通过对模型参数和合约特征的联合参数化，实现了训练一次、多参数组合即时定价。Guo等人\cite{guo2024icpinn}提出了改进约束PINN（ICPINN），通过输出变换将终值条件硬编码进网络结构，减少损失函数中的冲突项，加速收敛。

\subsection{三种独立PINN的复现与基准建立}\label{subsec:reproduction}

为提供可靠的对比基准，本文对上述三类代表性PINN方法进行了独立复现。

\textbf{BSM-PINN。}严格复现Dhiman和Hu\cite{dhiman2023pinn}的门控网络结构：浅层分支（2层$\times$64神经元，Tanh）与深层分支（4层$\times$64神经元，Tanh）通过可学习的sigmoid门加权融合，输入为归一化的$(S/S_{\max},\,t/T)$，输出为归一化期权价值$u=V/K$。相对于原论文，做了两处工程改动：（1）将上边界采样区间从单点$S_{\max}$扩展为$[0.8S_{\max},\,S_{\max}]$的均匀分布；（2）采用CosineAnnealingLR学习率调度（从$10^{-3}$衰减至$10^{-5}$）。以BSM解析解为参考，在$S\in[60,160]$、$t=0$处评估，MAE=0.016，RelMAE=0.6\%。

\textbf{CEV-PINN。}网络结构采用Dhiman和Hu\cite{dhiman2023pinn}的门控网络（hidden=64），将扩散系数替换为$\sigma^2 S^{2\beta}$，参数取$\sigma=0.20$、$\beta=0.5$。参考基准采用Schroder\cite{schroder1989computing}非中心卡方分布解析解：当$\beta=0.5$时，CEV欧式看涨期权价格为
\begin{equation}
  C = S_0\bigl[1 - F(\lambda K^{2(1-\beta)};\,d,\,\lambda S_0^{2(1-\beta)}e^{2r(1-\beta)T})\bigr]
    - Ke^{-rT} F\!\bigl(\lambda S_0^{2(1-\beta)}e^{2r(1-\beta)T};\,d-2,\,\lambda K^{2(1-\beta)}\bigr),
\end{equation}
其中$F(\cdot;\,d,\,\lambda)$为自由度$d=2+1/(1-\beta)$、非中心参数$\lambda=2r/[\sigma^2(1-\beta)(e^{2r(1-\beta)T}-1)]$的非中心卡方累积分布函数，可由\texttt{scipy.stats.ncx2}精确计算（精度$\sim10^{-15}$）。评估结果：MAE=0.074，RelMAE=1.69\%（仅统计$V^*>0.5$的区域）。

\textbf{Heston-ICPINN。}严格按照Tan和Zhang\cite{guo2024icpinn}的ICPINN方法复现，采用双网络结构：辅助网络（AuxNet，3层$\times$40神经元，ELU激活）预训练以满足Dirichlet边界条件，主网络（MainNet，8层$\times$40神经元，Tanh激活）学习PDE内部解。试验解形式为：
\begin{equation}\label{eq:icpinn_trial}
  U(S,v,\tau) = \phi(S) + S_n\tau_n \cdot \bigl[A(S_n,v_n,\tau_n) + \mathrm{NN}(S_n,v_n,\tau_n)\bigr],
\end{equation}
其中$\phi(S)=\max(S-K,0)$为解析嵌入的终值条件，$S_n\tau_n$在$\tau=0$和$S=0$处自动为零，从而将Dirichlet边界条件硬编码进网络结构。训练参数严格遵循论文设置：$K=100$，$r=0.1$，$\kappa=1$，$\theta=0.08$，$\xi=0.39$，$\rho=-0.93$，内部配点$N_r=10000$，边界配点$N_{bc}=500$，Adam优化器，初始学习率$10^{-3}$，StepLR调度（$\gamma=0.75$，每5000步），共训练20000轮。评估结果：MAE=0.35，RelMAE=3.9\%（仅统计实值期权）。

然而，ICPINN方法的收敛性对参数集高度敏感。当将参数替换为统一框架所需的参数集（$r=0.05$，$\kappa=2.0$，$\theta=0.04$，$\xi=0.3$，$\rho=-0.7$）时，训练无法收敛。根本原因在于Fichera退化边界条件（$v=0$处）：在$\kappa\theta$较小时，该边界条件的损失项主导整体损失（约为PDE残差的40倍），而aux\_net已被冻结，main\_net无法独立修正此偏差。为解决这一问题，本文在独立Heston-PINN中引入\textbf{数据锚点机制}：在$t=0$处预计算GL特征函数半解析价格作为监督信号，以权重$w_\text{data}=100$加入损失函数，将解锁定在正确分支上，从而在任意参数集下均能稳定收敛。

三种独立PINN的复现结果汇总于表~\ref{tab:reproduction}。

\begin{table}[htbp]
  \centering
  \caption{三种独立PINN复现结果（$S\in[60,160]$，$t=0$）}
  \label{tab:reproduction}
  \begin{tabular}{lcccc}
    \hline
    模型 & 参考基准 & MAE & RelMAE & 网络结构 \\
    \hline
    BSM-PINN      & 解析解           & 0.016 & 0.6\%  & 门控网络，2+4层$\times$64，Tanh \\
    CEV-PINN      & 非中心卡方解析解 & 0.074 & 1.69\% & 门控网络，2+4层$\times$64，Tanh \\
    Heston-ICPINN & 特征函数半解析解 & 0.35  & 3.9\%  & AuxNet(3$\times$40)+MainNet(8$\times$40) \\
    \hline
  \end{tabular}
\end{table}

上述三个独立PINN验证了各模型PDE的可解性，并为后续统一框架提供了精度基准。

\subsection{现有工作的局限性}

综合上述工作，现有PINN期权定价研究存在两方面不足：各研究针对单一模型，缺乏统一框架同时支持多种定价模型；用户需手动指定模型类型和参数，无法通过自然语言描述需求。本文针对上述不足，提出将三大模型统一为单一参数化PINN的方法，具体设计见第~\ref{chap:method}章。

\section{大语言模型辅助科学计算}\label{sec:llm_sci}

Liu等人\cite{liu2025pdeagent}提出了PDE-AGENT，一个工具链增强的多智能体框架，能够从自然语言描述出发，自动调用数值求解工具完成PDE的求解。该工作表明LLM具备理解PDE问题结构、选择合适求解策略的能力。在金融领域，Xiao等人\cite{xiao2024tradingagents}提出了TradingAgents框架，通过多个专业化LLM智能体的协作完成金融交易决策，验证了LLM在金融场景下的参数提取与任务路由能力。

现代LLM（如GPT-4、DeepSeek等）能够通过约束输出格式（如JSON模式）可靠地从自然语言中提取结构化参数。本文的LLM路由层设计较PDE-AGENT更为轻量：LLM仅负责参数提取和模型选择，不参与PDE求解过程，因此延迟更低（约500ms至2000ms的API调用）、可靠性更高（结构化JSON输出避免了自由文本解析的不确定性）。

\section{本章小结}\label{sec:related_summary}

本章梳理了BSM、CEV、Heston三大期权定价模型的PDE形式，回顾了PINN在期权定价中的代表性工作，并介绍了LLM辅助科学计算的进展。现有PINN工作验证了神经网络求解期权定价PDE的可行性，但缺乏统一的多模型框架和自然语言交互能力。本文在此基础上构建了LLM与PINN相结合的统一期权定价系统，具体方法见第~\ref{chap:method}章。
